% =====================================================================
%  12. Organisation du projet et diagramme de Gantt
% =====================================================================
\clearpage
\section{Organisation du projet et diagramme de Gantt}
\label{sec:organisation}

\subsection{Méthode de travail}

Le travail a été organisé autour d'un dépôt Git commun (\cref{tab:depot}) :
une branche par tâche, des modèles d'\emph{issue} et de \emph{pull request}, une
spécification vivante (\fichier{specs/spec.md}) consignant chaque décision
scientifique, des rapports intermédiaires datés et une suite de tests ciblés.
Le fichier \fichier{specs/time\_log.csv} est prévu pour enregistrer le temps
passé par membre et par séance.

\subsection{Chronologie et diagramme de Gantt}

La \cref{fig:gantt} retrace les grandes phases du projet. Les dates proviennent
de l'historique Git et des dates portées par les rapports intermédiaires : elles
indiquent \emph{quand} chaque livrable a été produit ou versionné, pas la durée
de travail effective.

\begin{figure}[htbp]
\centering
\begin{ganttchart}[
    hgrid, vgrid={*1{draw=black!12}},
    time slot format=isodate,
    x unit=0.98cm, y unit chart=0.58cm, y unit title=0.55cm,
    title height=1,
    title label font=\scriptsize\bfseries,
    bar label font=\scriptsize,
    group label font=\scriptsize\bfseries,
    milestone label font=\scriptsize\itshape,
    bar/.append style={fill=csbordeaux!70, draw=csbordeaux},
    bar height=0.55,
    group/.append style={fill=csgris, draw=csgris},
    group height=0.25,
    milestone/.append style={fill=orangeviz, draw=orangeviz},
    bar label node/.append style={align=right},
  ]{2026-10-02}{2026-10-11}
  \gantttitle{Octobre 2026}{10} \\
  \gantttitlecalendar{day} \\
  \ganttgroup{Cadrage et données}{2026-10-02}{2026-10-07} \\
  \ganttbar{Dépôt, énoncé, téléchargement}{2026-10-02}{2026-10-03} \\
  \ganttbar{Audit exploratoire (EDA)}{2026-10-03}{2026-10-07} \\
  \ganttbar{Essais de prétraitement, ACP}{2026-10-05}{2026-10-05} \\
  \ganttgroup{Préparation et protocole}{2026-10-08}{2026-10-09} \\
  \ganttbar{Audit et migration (notebook 02)}{2026-10-08}{2026-10-08} \\
  \ganttbar{Cibles et plan d'expériences}{2026-10-09}{2026-10-09} \\
  \ganttgroup{Modélisation et analyse}{2026-10-10}{2026-10-10} \\
  \ganttbar{Supervisé et diagnostics (03, 04)}{2026-10-10}{2026-10-10} \\
  \ganttbar{CV aléatoire, exports (05, 06)}{2026-10-10}{2026-10-10} \\
  \ganttbar{COREG (07)}{2026-10-10}{2026-10-10} \\
  \ganttgroup{Rédaction}{2026-10-10}{2026-10-11} \\
  \ganttbar{Synthèses et rapports}{2026-10-10}{2026-10-11} \\
  \ganttmilestone{Rendu final (date à compléter)}{2026-10-11}
\end{ganttchart}
\caption{Diagramme de Gantt des phases du projet, d'après l'historique Git et les
dates des rapports intermédiaires. Les responsables et les heures réelles sont
à renseigner dans le \cref{tab:taches} ; la date de rendu est indicative.}
\label{fig:gantt}
\end{figure}

\subsection{Répartition des tâches et temps passé}

L'énoncé demande d'indiquer, pour chaque tâche, le temps qu'elle a pris et le
membre de l'équipe qui s'en est chargé. Ces informations ne sont pas déductibles
des fichiers du dépôt (le journal de temps est vide) : elles sont laissées
\emph{à compléter par l'équipe}, sans être inventées. Les durées de calcul des
notebooks (quelques minutes chacun) ne sont pas des durées de travail humain.

\begin{table}[htbp]
\centering
\small
\caption{Répartition des tâches. Colonnes \og Membre(s) \fg{} et \og Heures \fg{} à
compléter par l'équipe à partir de \fichier{specs/time\_log.csv}.}
\label{tab:taches}
\begin{tabularx}{\textwidth}{@{}L l l l@{}}
\toprule
Tâche & Livrables & Membre(s) & Heures \\
\midrule
Cadrage, dépôt, téléchargement des données & \fichier{scripts/}, \fichier{specs/} & \acompleter{} & \acompleter{} \\
Analyse exploratoire et audit des mesures & notebook 01, rapport 01 & \acompleter{} & \acompleter{} \\
Prétraitement, groupes, cohorte, plis & notebook 02, rapports 02--08 & \acompleter{} & \acompleter{} \\
Plan d'expériences, choix des cibles & rapport 09, \fichier{spec.md} & \acompleter{} & \acompleter{} \\
Entraînement supervisé & notebooks 03, 05, 06 & \acompleter{} & \acompleter{} \\
Diagnostics et explicabilité & notebook 04, rapport 11 & \acompleter{} & \acompleter{} \\
Bibliographie et COREG & notebook 07, rapport 14 & \acompleter{} & \acompleter{} \\
Tests et reproductibilité & \fichier{tests/}, manifestes & \acompleter{} & \acompleter{} \\
Rédaction et relecture croisée & rapports 15, 16, présent rapport & \acompleter{} & \acompleter{} \\
\midrule
\textbf{Total} & & & \acompleter{} \\
\bottomrule
\end{tabularx}
\end{table}

\begin{aretenir}
Le projet s'est déroulé du 2 au 10 octobre 2026 pour sa partie technique, en
quatre phases : cadrage et données, préparation et protocole, modélisation et
analyse, rédaction. La répartition des tâches et les heures réelles restent à
renseigner par l'équipe.
\end{aretenir}
